\subsubsection{A wild cover above one low-degree positive-genus curve}

\begin{example}[Refuted: universal tameness above a single smooth exceptional curve]
\label{D:wild-genus2-single-curve}
Let $k$ be algebraically closed of characteristic two.  There is
a normal affine surface $S$, with only one singular point $o$,
whose resolution has exceptional fibre one smooth genus-two
curve $C$ with $C^2=-1$, and a connected finite Galois cover
$T\to S$ with group $\mathbb Z/2\mathbb Z$ which is \'etale
over $S\setminus\{o\}$ and wild at $C$.  At its unique
prolongation the ramification index is one, the residue
extension is purely inseparable of degree two, and the
different exponent is one.

We can take the following explicit hypersurfaces:
\[
 \begin{aligned}
 S&:\quad F(a,b,t)=b^2+(t^5+t^2a^2)b+a^5=0,\\
 T&:\quad H(c,d,t)=d^2+(c^4+t^3+t^4c^2)d
            +tc^5+t^2c^7+t^4c^6+t^6c^5=0,
 \end{aligned}
\]
with map
\[
                 a=c^2+t^2c,\qquad b=t^2d+c^5.
 \tag{M}
\]
In particular the hypothesis of formal linearity in
Proposition~\ref{D:graded-cone-tame} cannot be discarded,
even for a Gorenstein singularity with one smooth exceptional
curve.  This example is a ramification obstruction; it is not
a counterexample to either finite-generation question D.1 or
D.2(ii).
\end{example}

\begin{proof}
All polynomial calculations in this proof take place in
characteristic two.

\par\smallskip\noindent\textbf{The target and a resolution.}\quad
The polynomial $F$ is irreducible.  Indeed, over $k(t)(a)$ put
$h=t^5+t^2a^2$ and divide its quadratic equation by $h^2$.
A root would give an Artin--Schreier expression
$a^5/h^2=w^2+w$.  The left-hand side has a pole of order one
at $a=\infty$.  Every pole of a rational function $w^2+w$
has even order, a contradiction.  The derivatives are
\[
             F_a=a^4,\qquad F_b=t^2(a^2+t^3),
             \qquad F_t=bt^4.
\]
Their common zero on $S$ is only $o=(0,0,0)$.  Thus $S$ is
an integral hypersurface, regular in codimension one and
Cohen--Macaulay, hence normal by Serre's criterion
\cite[Lemma 10.157.4, tag 031S]{StacksBWildCurve}.

We give an explicit proper resolution.  On $k[a,b,t]$ use
weights $(2,5,1)$, let $I_m$ be the ideal spanned by monomials
of weight at least $m$, and form
\[
 \mathcal R=\bigoplus_{m\geq0}I_m s^m.
\]
This is a finite type graded $k[a,b,t]$-algebra: it is generated
by $as,as^2$, by $bs,\ldots,bs^5$, and by $ts$.
The morphism $\operatorname{Proj}\mathcal R\to\mathbb A^3$
is proper, by the relative-Proj theorem
\cite[Lemma 31.31.4, tag 07ZW]{StacksBWildCurve}, and is an
isomorphism away from the origin.  Let $Y$ be the strict
transform of $S$.  Its reduced exceptional fibre is the
weighted projective curve
\[
               b^2+t^5b+a^5=0
                    \quad\hbox{in }\mathbb P(2,5,1).
 \tag{E}
\]
We verify smoothness using charts, so no smoothness assertion
about an ambient weighted projective space is needed.

On the $t$ chart put $a=t^2x$ and $b=t^5y$.  Dividing $F$ by
$t^{10}$ gives
\[
                    y^2+(1+tx^2)y+x^5=0.
 \tag{Y0}
\]
This chart is smooth: its partial derivatives with respect to
$x,y$ are $x^4$ and $1+tx^2$, which never vanish together.
It contains all of the exceptional curve except its point at
infinity.  Near that point both weighted coordinates $a,b$
are nonzero.  The degree-zero ring of the weighted blowup
chart has coordinates
\[
 \rho=\frac b{a^2},\qquad
 \xi=\frac{a^5}{b^2}\in\mathcal O^\times,
 \qquad \tau=\frac{ta^2}{b}.
\]
The inverse monomial formulas are
\[
                  a=\rho^2\xi,
       \qquad b=\rho^5\xi^2,
       \qquad t=\rho\tau.
\]
Their exponent matrix is unimodular.  The weighted inequalities
give exactly the ring $k[\rho,\tau,\xi,\xi^{-1}]$ on this
ambient chart.  Dividing the substituted equation by
$\rho^{10}\xi^2$ gives the strict transform
\[
                  \xi^2(1+\xi+\rho\tau^2)+\tau^5=0.
 \tag{Yinf}
\]
Its derivative with respect to $\xi$ is $\xi^2$, a unit.
The unique exceptional point outside (Y0) is
$(\rho,\tau,\xi)=(0,0,1)$.  The two charts cover the
exceptional fibre, so $Y$ is smooth there.  Away from that
fibre it equals the already smooth $S\setminus\{o\}$.
Thus $Y\to S$ is a proper resolution with only the curve
$C$ of (E) exceptional.

For completeness, $C$ is the smooth proper completion of
$y^2+y=x^5$.  At infinity use $u=1/x$ and $v=y/x^3$; its
equation becomes $v^2+u^3v=u$, smooth at $(u,v)=(0,0)$.
The finite double cover $C\to\mathbb P^1_x$ has
\[
                 \pi_*\mathcal O_C
                    =\mathcal O\oplus\mathcal O(-3),
\]
as follows directly from the bases $1,y$ and $1,v$ and the
transition $y=x^3v$.  Hence $\chi(\mathcal O_C)=1-2=-1$
and $g(C)=2$.  It is integral because $x^5$ cannot be
$w^2+w$ in $k(x)$, by the same odd-pole argument.
On (Yinf), the divisor of the original function $t$ is
locally $\{\rho=0\}+\{\tau=0\}$.  The second curve meets
$C$ transversally once, and there are no other intersections
with the strict transform of $t=0$.  Intersecting the
principal divisor of $t$ with the proper curve $C$ yields
\[
                         0=C^2+1.
\]
Thus $C^2=-1$.

\par\smallskip\noindent\textbf{The connected quadratic extension.}\quad
On the $t$ chart the required extension of function fields is
\[
                            z^2+tz=x.
 \tag{AS}
\]
We verify its nontriviality, including the global gluing which
explains this choice.  Over $k[t]$, the two equations
\[
 y^2+(1+tx^2)y=x^5,
 \qquad v^2+(u^3+tu)v=u
\]
with $u=1/x$, $v=y/x^3$ define a smooth proper genus-two
curve family $\mathcal C\to\mathbb A^1_t$.  Smoothness on
the first chart was checked above; on the second chart the
only new point has $u=v=0$, where the $u$ derivative is one.
It is finite flat of degree two over $\mathbb P^1\times
\mathbb A^1$, so proper.  Again the direct-image module is
$\mathcal O\oplus\mathcal O(-3)$.

The \v Cech class
\[
                             h=\frac y{x^2}
\]
is a nonzero basis vector of $H^1(\mathcal C_t,\mathcal O)$,
including over $k(t)$: the two basis classes are $y/x$ and
$y/x^2$.  The equation of the family gives
\[
                         h^2+th=x+\frac y{x^4}.
\]
Consequently the two covers
\[
 z_0^2+tz_0=x,
       \qquad z_\infty^2+tz_\infty=uv
\]
glue by $z_\infty=z_0+h$.  They form a torsor under the
finite flat group scheme $\ker(F-t:\mathbb G_a\to\mathbb G_a)$.
On $t\ne0$ this is an \'etale group of order two, and the
torsor maps to the nonzero class $h$ under the map to
$H^1(\mathcal O)$ from its defining additive exact sequence.
Thus its generic fibre is nontrivial and connected.  This
proves that (AS) defines a separable quadratic field extension.

Put $c=tz$ and $d=(b+c^5)/t^2$ in that extension.  Direct
substitution gives
\[
          F(c^2+t^2c,t^2d+c^5,t)=t^4H(c,d,t).
 \tag{I}
\]
The algebra defining $T$ is finite over that of $S$: $c$
satisfies the monic equation $c^2+t^2c=a$, and $d$ satisfies
the displayed monic equation $H=0$ over $k[c,t]$.
On $t\ne0$ it is exactly the quadratic cover (AS).  Therefore
it is integral: its localization at $t$ is a domain and the
monic hypersurface equation has no component supported on
$t=0$.

The scheme $T$ is smooth away from the origin.  On $t\ne0$
this follows from the \'etale description over the smooth
target.  On $t=0$ its equation is $d(d+c^4)=0$ and its
$d$ derivative is $c^4$, so only $(c,d,t)=(0,0,0)$ can
be singular.  Being an integral hypersurface, $T$ is normal
by the same $(R_1)+(S_2)$ argument.  Thus it is the finite
normalization of $S$ in the quadratic extension.
The nontrivial deck transformation is explicitly
\[
       c\longmapsto c+t^2,
       \qquad d\longmapsto d+c^4+t^6c+t^8.
\]
It preserves (M) and $H$, and its square is the identity.

\par\smallskip\noindent\textbf{\'Etaleness away from the singular point.}\quad
On $t\ne0$ the derivative of $c^2+t^2c-a$ with respect to
$c$ is $t^2$, so the cover is \'etale.  If $t=0$ and the
image is not $o$, then $a\ne0$, and hence $c\ne0$.
Both surfaces are smooth at such points.  The target has
local \'etale coordinates $(b,t)$, since $F_a=a^4$ is a
unit.  The source has local \'etale coordinates $(c,t)$,
since $H_d=c^4$ at $t=0$.  In these coordinates
\[
                  db=c^4\,dc,
                         \qquad dt=dt
                    \quad\hbox{at }t=0.
\]
The differential is invertible.  Equivalently, the finite map
is flat by miracle flatness and its relative differentials
vanish; see \cite[Lemma 10.128.1, tag 00R4]{StacksBWildCurve}.
It is therefore \'etale at all these points.  The fibre above
$o$ is only the origin, so this proves the claimed
quasi-\'etale property.

\par\smallskip\noindent\textbf{The exceptional valuation.}\quad
At the generic point of $C$ in (Y0), $t$ is a uniformizer
and the residue field is
\[
                 k(C)=k(x,y),\qquad y^2+y=x^5.
\]
The element $x$ is not a square in $k(C)$: the separable map
to $k(x)$ implies $dx\ne0$ (or the displayed curve equation
gives the same differential calculation).  The reduction of
(AS) is therefore the irreducible purely inseparable equation
\[
                              \bar z^2=x.
\]
The algebra obtained by adjoining $z$ to the boundary DVR is
local, finite free of rank two, and has maximal ideal generated
by $t$.  It is consequently a regular DVR and is already the
integral closure.  We obtain
\[
                  e=1,
           \qquad [k(C)(\bar z):k(C)]_{\mathrm{insep}}=2.
\]
The deck action
$z\mapsto z+t$ is trivial on that residue field.  Thus all
of $\mathbb Z/2\mathbb Z$ is inertia and is wild.  The monogenic
different is generated by the derivative $t$ of (AS), so its
exponent is $\delta=1$; see
\cite[Example 49.12.5, tag 0BWB]{StacksBWildCurve}.

The target is Gorenstein because it is a hypersurface.
Adjunction gives $K_Y\cdot C=2g-2-C^2=3$; the discrepancy
of the sole exceptional curve is therefore $-3$.  In
particular the example lies outside the log-canonical case.
It also lies outside the formal-cone case, since that case
has tame exceptional ramification by
Proposition~\ref{D:graded-cone-tame}.  The inequality of
Proposition~\ref{D:single-curve-residue-tameness} is sharp
here: its two sides equal $2$ with $g=2,d=1,n=2,e=1$
and inseparable residue degree two.

The companion program \texttt{check\_wild\_genus2\_curve.py}
performs 26 exact assertions over $\mathbb F_2$, including
(I), the two Jacobians, the involution, both resolution
charts, and the \v Cech identities.  Its recorded runtime
is about $0.15$ seconds.  The proofs of properness, genus,
connectedness, and ramification above are separate from
these finite polynomial certificates.
\end{proof}
